\documentclass[11pt]{amsart}

\usepackage[T1]{fontenc}
\usepackage[utf8]{inputenc}
\usepackage{lmodern}
\usepackage{amssymb}
\usepackage{mathtools}
\usepackage[margin=1in]{geometry}
\usepackage[hidelinks]{hyperref}

\setlength{\emergencystretch}{3em}

\newcommand{\lean}[1]{\texttt{\detokenize{#1}}}
\newcommand{\fn}[1]{\mathrm{#1}}
\newcommand{\st}[2]{[#1]_{#2}}
\newcommand{\leanref}[1]{\par\nopagebreak\smallskip
  {\small\raggedright\noindent Lean: #1\par}\medskip}

\newtheorem{theorem}{Theorem}
\theoremstyle{definition}
\newtheorem{definition}{Definition}
\newtheorem{example}{Example}

\title{Irreducible moduli: a formally verified residue number system over arbitrary moduli}
\author{Artur Flamandzki}
\date{}

\begin{document}

\begin{abstract}
We define a residue number system over arbitrary moduli, composite and not pairwise coprime
moduli included. Residues are computed by counting, greatest common divisors by Euclid's
algorithm on residues, and the moduli are combined into a mixed-radix system in which each
modulus $m$ receives the radix $m / \gcd(W, m)$, where $W$ is the product of the radices before
it. No definition refers to primes, to factorization, or to a number assembled from residues.
We prove that the product of the radices is the least common multiple of the moduli, in any
order; that on the consecutive moduli $2, \dots, n-1$ the modulus $n$ receives the radix $p$ when
$n = p^k$ and $1$ otherwise, so it receives its whole value exactly when $n$ is prime, and the
logarithm of the radix is the von Mangoldt function. Numbers given only by their residues can
be compared, tested for primality, added, subtracted, multiplied, raised to powers and divided
exactly, with overflow and underflow detected; and the residue computed by counting is also
read from a comb of hyperbolic-secant pulses from scale $3$ on. The work states several
independent topics in one formal frame: remainders, Euclid's algorithm, the least common
multiple of consecutive integers and the von Mangoldt function, primes in intervals of
consecutive integers, and residue number system arithmetic over arbitrary moduli. Every
statement in this paper is a theorem of the accompanying Lean~4 / Mathlib formalization, which
builds with no \lean{sorry} and uses only the axioms \lean{propext}, \lean{Quot.sound} and
\lean{Classical.choice}.
\end{abstract}

\maketitle

\section{Introduction}

A residue number system represents a number by its residues modulo a list of moduli. This
paper defines such a system over arbitrary moduli, using counting modulo a modulus as the only
primitive, and proves what it computes.

The definitions (Section~\ref{sec:definitions}) use only moduli and residues below a modulus.
Primes, factorizations and numbers assembled from residues do not occur in them. They occur
only in the theorems (Section~\ref{sec:results}), which identify the computed quantities with
the usual arithmetic ones: residues computed by counting with remainders, Euclid's algorithm on
residues with greatest common divisors, the product of the radices with the least common
multiple of the moduli, and the moduli that receive their whole value as radix after the
consecutive moduli below them with the primes.

\subsection*{Purpose and contributions}
The aim is one frame in which several topics that are usually treated separately are stated
with the same few notions and verified by the same proof checker. Within that frame, the work
contributes three points that the author did not find in the literature:
\begin{enumerate}
  \item[(C1)] A machine-checked residue number system with its non-modular operations:
    magnitude comparison, detection of overflow and underflow, exact division, and a primality
    test, all on the residues of a number (T5, T5c, T6, T8). The author found formal proofs of
    the Chinese remainder theorem, but no formal verification of comparison or mixed-radix
    conversion in a residue number system.
  \item[(C2)] Arbitrary moduli. The moduli need not be pairwise coprime: each modulus $m$
    receives the radix $m / \gcd(W, m)$, only the part of $m$ that the moduli before it do not
    already cover (T3), and comparison works on these radices (T5). The comparison methods the
    author found assume pairwise coprime moduli.
  \item[(C3)] No number is assembled from residues. Every computation goes through residues;
    the only number formed from residues is the residue of a block modulus, which is below that
    modulus (Section~\ref{sec:blocks}).
\end{enumerate}
In the same frame the work proves several known results (T1, T2, T4, T7) from the definitions
rather than restating them, and adds the reading of residues from a sech comb (T1w).

The comparison with earlier work reflects the literature known to the author at the time of
writing; it states what the author did not find, not that no such work exists.

\subsection*{Related work}
Theorems T2--T4 are classical. The identity
$\fn{lcm}(1, \dots, n) = e^{\psi(n)}$, with the jump by $p$ exactly at the powers of $p$, is
formalized in Isabelle by Eberl~\cite{eberl-afp} and is part of Mathlib~\cite{mathlib} as
\lean{Chebyshev.psi_eq_log_lcmUpto}; Mathlib also contains the Chinese remainder theorem for
moduli that are not coprime (\lean{Nat.chineseRemainder'}). Here these facts are derived from
definitions that do not mention them.

Residue number systems go back to Garner~\cite{garner}. Addition and multiplication act on each
residue separately, while comparison, overflow detection and division need a positional
representation of the number, usually through mixed-radix conversion; a recent method and an
account of the earlier ones are given in~\cite{didier}. Moduli that are not coprime are used in
error-correcting residue codes~\cite{sundaram}. For the reading of residues from a sech comb
(T1w) the author found no earlier work.

Every theorem below is stated with the name of the Lean theorem that proves it. Hypotheses are
given in full; nothing here is weaker or stronger than the Lean statement. The formalization
and the build instructions are described in Section~\ref{sec:formalization}.

\subsection*{Notation}
All numbers are natural numbers. We write $x \bmod 0 = x$ (the Lean convention), and
$\lfloor x/y \rfloor$ for integer division, with $\lfloor x/0 \rfloor = 0$. A list of moduli is
written $R = (r_1, \dots, r_k)$; ``positive moduli'' means $r_i \ge 1$ for every $i$. We write
$W(R) = \fn{lcm}(r_1, \dots, r_k)$, with $W = 1$ for the empty list; the residue vector of $x$
on $R$ is $\st{x}{R} = (x \bmod r_1, \dots, x \bmod r_k)$; and $\prod R = r_1 \cdots r_k$. The
list $R$ followed by $m_1, \dots, m_j$ is written $(R, m_1, \dots, m_j)$. $\Lambda$ is the von
Mangoldt function.

\section{Definitions}\label{sec:definitions}

\subsection{Residues by counting and Euclid's algorithm}

\begin{definition}
The residue of $t$ modulo $d$ is computed by counting alone: add $1$ per step and reset to $0$
on reaching $d$,
\[
  \fn{countMod}(d, 0) = 0, \qquad
  \fn{countMod}(d, t + 1) =
  \begin{cases}
    0 & \text{if } \fn{countMod}(d, t) + 1 = d, \\
    \fn{countMod}(d, t) + 1 & \text{otherwise,}
  \end{cases}
\]
and $\fn{resid}(a, b) = \fn{countMod}(b, a)$.
\end{definition}
\leanref{\lean{countMod}, \lean{resid}}

\begin{definition}
Euclid's algorithm on residues:
\[
  \fn{euclid}(a, 0) = a, \qquad
  \fn{euclid}(a, b) = \fn{euclid}\big(b, \fn{resid}(a, b)\big) \quad (b \ge 1).
\]
In Lean the recursion carries a fuel bound of $b + 1$ steps, which is always enough.
\end{definition}
\leanref{\lean{euclid}}

\subsection{The mixed-radix system}

\begin{definition}
The residue modulo $c$ of the product of a list of radices is computed factor by factor,
without forming the product:
\[
  \fn{prodMod}((), c) = 1 \bmod c, \qquad
  \fn{prodMod}((r, R'), c) = \big(\fn{resid}(r, c) \cdot \fn{prodMod}(R', c)\big) \bmod c.
\]
The gcd of $m$ with that product is
\[
  \fn{gcdProd}(R, m) = \fn{euclid}\big(\fn{prodMod}(R, m), m\big).
\]
\end{definition}
\leanref{\lean{prodMod}, \lean{gcdProd}}

\begin{definition}
The mixed-radix system of the moduli $m_1, \dots, m_k$ in this order has the radices
$L_k = \fn{radices}(m_1, \dots, m_k)$, where $L_0 = ()$ and
\[
  L_{j+1} = \Big(L_j,\ \frac{m_{j+1}}{\fn{gcdProd}(L_j, m_{j+1})}\Big).
\]
The radix modulus $m$ receives after the moduli $B$ is
\[
  \fn{newRadix}(B, m) = \frac{m}{\fn{gcdProd}(\fn{radices}(B), m)},
\]
and $m$ is irreducible after $B$, $\fn{IrreducibleAfter}(B, m)$, when
$\fn{newRadix}(B, m) = m$: it receives its whole value as radix. Finally
$\fn{initialModuli}(n) = (2, 3, \dots, n - 1)$.
\end{definition}
\leanref{\lean{radices}, \lean{newRadix}, \lean{IrreducibleAfter}, \lean{initialModuli}}

\begin{example}
Checked by the Lean kernel (\lean{decide}):
\[
  \fn{radices}(2, 3, 4, 5, 6, 7, 8, 9, 10) = (2, 3, 2, 5, 1, 7, 2, 3, 1),
\]
\[
  \fn{radices}(6, 35) = (6, 35), \qquad \fn{radices}(4, 4) = (4, 1).
\]
\end{example}

\subsection{Mixed-radix representations and comparison}

A mixed-radix representation is a list of pairs of a digit and a radix, least significant
first: $D = \big((d_1, r_1), \dots, (d_k, r_k)\big)$.

\begin{definition}
The residue modulo $c$ of the value of $D$ is computed by Horner's rule from the most
significant digit, with every radix reduced modulo $c$:
\begin{align*}
  \fn{evalMod}((), c) &= 0, \\
  \fn{evalMod}\big(((d, r), D'), c\big) &= \big(d + \fn{resid}(r, c) \cdot \fn{evalMod}(D', c)\big) \bmod c.
\end{align*}
The value of $D$ is $\fn{eval}(D) = \sum_{i=1}^{k} d_i \prod_{j < i} r_j$. It is formed only as
the residue of a block modulus (Section~\ref{sec:blocks}), where it is below that modulus.
\end{definition}
\leanref{\lean{evalMod}, \lean{eval}}

\begin{definition}
$\fn{invMod}(a, b)$ is the least $x \ge 1$ with $x a \equiv 1 \pmod b$, found by repeated
addition of $a$ modulo $b$ (at most $b$ additions; $\fn{invMod}(a, 1) = 0$).
$\fn{toMixedRadix}(R, S)$ converts the residues $S$ on the moduli $R$ to a mixed-radix
representation, modulus by modulus (Garner's step on arbitrary moduli): modulus $m$ with
residue $s$, appended to the representation $D$ built so far, with
$g = \fn{gcdProd}(\text{radices of } D, m)$ and $r = m / g$, takes
$y = (s - \fn{evalMod}(D, m)) \bmod m$; if $g \nmid y$ the residues are incompatible and the
result is $\fn{none}$; otherwise the new radix is $r$ and the new digit is
$(y / g) \cdot \fn{invMod}\big((\fn{prodMod}(\text{radices of } D, m) / g) \bmod r,\ r\big) \bmod r$
(digit $0$ when $r = 1$). $\fn{compareRes}(R, S, S')$ converts both residue vectors and
compares the two representations digit by digit from the most significant.
\end{definition}
\leanref{\lean{invMod}, \lean{toMixedRadix}, \lean{compareRes}}

\subsection{Blocks}\label{sec:blocks}

A block is a list of moduli used as one modulus: its modulus is the lcm of its members, and its
residue is the value of the mixed-radix representation of the member residues.

\begin{definition}
A block of members $M$ with residues $S$ is the modulus and residue
\[
  \fn{blockModulus}(M, S) = \big(\textstyle\prod \text{radices of } D,\ \fn{eval}(D)\big),
  \qquad D = \fn{toMixedRadix}(M, S),
\]
and $\fn{blockModuli}$ applies $\fn{blockModulus}$ to each block of a list. Base extension,
$\fn{baseExtend}(R, S, T)$, takes $D = \fn{toMixedRadix}(R, S)$ and gives $\fn{evalMod}(D, c)$
for each modulus $c$ in $T$.
\end{definition}
\leanref{\lean{blockModulus}, \lean{blockModuli}, \lean{baseExtend}}

\begin{definition}
A number in block representation ($\fn{BlockRep}$) is given by the member residues of its
blocks and, for its quotient by the lcm of those blocks, the member residues of its quotient
blocks. $\fn{flatten}$ turns it into one list of moduli: the block moduli followed by the
quotient block moduli, where a quotient block of modulus $c$ with residue $k$ receives the
residue
$\big(k \cdot \fn{prodMod}(\text{radices of the blocks}, c) + \fn{evalMod}(\text{representation of the blocks}, c)\big) \bmod c$.
$\fn{compareBlocks}(x, y)$ flattens both numbers, extends both to the moduli of $x$ followed by
the moduli of $y$ not among them, and compares them there with $\fn{compareRes}$.

For statements only, $\fn{BlockRep.of}(n, D, K)$ writes $n$ in blocks $D$ and quotient blocks
$K$: the members $m$ of $D$ receive the residues $n \bmod m$, those of $K$ the residues
$\lfloor n / L \rfloor \bmod m$, where $L$ is the lcm of all members of $D$.
\end{definition}
\leanref{\lean{BlockRep}, \lean{flatten}, \lean{compareBlocks}, \lean{BlockRep.of}}

\subsection{Primality on residues}

\begin{definition}
A number is given by its residues $S$ on moduli $R$. Trial divisors $m = 2, 3, \dots$, at most
$\fn{fuel}$ of them, are compared with it through their own residue vectors
$(\fn{resid}(m, c))_{c \in R}$ and $(\fn{resid}(m, c)^2 \bmod c)_{c \in R}$, while the residue
of the number modulo $m$ is computed with $\fn{evalMod}$ from its mixed-radix representation.
$\fn{isPrimeRes}(\fn{fuel}, R, S)$ is $\fn{true}$ when $m$ equals the number or $m \cdot m$
exceeds it, $\fn{false}$ when the number is below $m$ or divisible by $m$, and $\fn{none}$ when
the fuel runs out or the residues are incompatible.
\end{definition}
\leanref{\lean{isPrimeRes}}

\subsection{Atoms of an interval}

\begin{definition}
The differences of the interval $lo, \dots, hi$ are
\[
  \fn{Diff}(lo, hi) = \{\, k \ge 2 : \exists a:\ lo \le a,\ k \mid a,\ a + k \le hi \,\}.
\]
\end{definition}
\leanref{\lean{IntervalDifference}}

\begin{definition}
$\fn{firstMultipleFrom}(lo, k)$ is the least multiple of $k$ that is at least $lo$, computed
from one residue: $lo$ if $\fn{resid}(lo, k) = 0$, and $lo + k - \fn{resid}(lo, k)$ otherwise.
The atoms of the interval $lo, \dots, hi$, $\fn{intervalAtoms}(lo, hi)$, are the differences
$k = 2, \dots, hi - lo$, taken in order, that are kept when the interval contains $a$ and
$a + k$ with $a = \fn{firstMultipleFrom}(lo, k)$, and $k$ is divisible by none of the atoms
kept before. The definition does not refer to divisibility order; Theorem T7 shows that the
atoms are the minimal elements of $\fn{Diff}(lo, hi)$ under divisibility, that is, the atoms
of $\fn{Diff}(lo, hi)$ ordered by divisibility.
\end{definition}
\leanref{\lean{firstMultipleFrom}, \lean{intervalAtoms}}

\subsection{Arithmetic on residues}

\begin{definition}
For each modulus $m$ with residues $a$ and $b$: $\fn{addRes}$ gives $(a + b) \bmod m$,
$\fn{subRes}$ gives $(a + (m - b)) \bmod m$, $\fn{mulRes}$ gives $a b \bmod m$; $\fn{powRes}$
performs $e$ products starting from the residues of $1$. A sum overflows, $\fn{addOverflow}$,
when its residue vector compares below that of the first term; a difference underflows,
$\fn{subUnderflow}$, when the first term compares below the second.
\end{definition}
\leanref{\lean{addRes}, \lean{subRes}, \lean{mulRes}, \lean{powRes}, \lean{addOverflow},
\lean{subUnderflow}}

\begin{definition}
For a divisor $t$ each modulus $m$ is replaced by $m / \fn{euclid}(m, t)$; we write
$R/t = \fn{divModuli}(R, t)$ for the resulting list. On modulus $m$ with
$g = \fn{euclid}(m, t)$ and $r = m / g$, $\fn{exactDivRes}$ maps the residue $s$ to
$(s / g) \cdot \fn{invMod}\big((t / g) \bmod r,\ r\big) \bmod r$. $\fn{divModRes}(R, S, T)$
computes the remainder $\rho$ modulo $T$ with $\fn{evalMod}$, subtracts $\rho$ modulus by
modulus and divides exactly by $T$. $\fn{addressRes}$ gives the address $(i, k)$ with the
offsets counted $1, \dots, T$: remainder $0$ becomes offset $T$, and otherwise the quotient is
increased by one.
\end{definition}
\leanref{\lean{divModuli}, \lean{exactDivRes}, \lean{divModRes}, \lean{addressRes}}

\subsection{The sech comb}

\begin{definition}
Modulus $d$ is represented by a sum of $\fn{sech}$ pulses of width $1/s$ centred at the
multiples of $d$, where $s$ is the scale:
\[
  \fn{sech}(x) = \frac{1}{\cosh(x)}, \qquad
  C_{s,d}(t) = \sum_{i=-1}^{2} \fn{sech}\Big(s \big(t - d (\lfloor t/d \rfloor + i)\big)\Big).
\]
A step $t$ is a peak for $d$ when $C_{s,d}(t) > \frac12$. $\fn{lastPeak}_s(d, t)$ is the last
peak for $d$ among the steps $1, \dots, t$ ($0$ if there is none), and
$\fn{combResid}_s(a, b) = a - \fn{lastPeak}_s(b, a)$ reads the residue of $a$ modulo $b$
backwards from $a$ to the last peak for $b$.
\end{definition}
\leanref{\lean{sech}, \lean{sechComb}, \lean{lastPeak}, \lean{combResid}}

\section{Results}\label{sec:results}

\subsection*{T1 Counting}
\begin{theorem}
For every $d$ and $t$, $d = 0$ included, $\fn{countMod}(d, t) = t \bmod d$. Hence
$\fn{resid}(a, b) = a \bmod b$.
\end{theorem}
\leanref{\lean{countMod_eq_mod}, \lean{resid_eq_mod}}

\subsection*{T1w Sech comb}
\begin{theorem}
For scale $s \ge 3$, every modulus $d \ge 1$ and every step $t$,
\[
  C_{s,d}(t) > \tfrac12 \iff d \mid t .
\]
Hence for $s \ge 3$ and $b \ge 1$, $\fn{combResid}_s(a, b) = \fn{resid}(a, b)$. At $s = 2$
the reading fails: $\fn{combResid}_2(3, 2) = 0 \ne 1 = \fn{resid}(3, 2)$.
\end{theorem}
\leanref{\lean{one_half_lt_sechComb_iff}, \lean{combResid_eq}, \lean{combResid_two_ne}}

\subsection*{T2 Euclid}
\begin{theorem}
$\fn{euclid}(a, b) = \gcd(a, b)$.
\end{theorem}
\leanref{\lean{euclid_eq_gcd}}

\subsection*{T3 Mixed-radix system}
\begin{theorem}
For any list $R$ and any $c$, $m$:
\[
  \fn{prodMod}(R, c) = \textstyle\prod R \bmod c, \qquad
  \fn{gcdProd}(R, m) = \gcd\big(\textstyle\prod R, m\big),
\]
\[
  \fn{radices}(R, m) = \Big(\fn{radices}(R),\ \frac{m}{\gcd\big(\prod \fn{radices}(R), m\big)}\Big).
\]
\end{theorem}
\leanref{\lean{prodMod_eq}, \lean{gcdProd_eq}, \lean{radices_append}}

\begin{theorem}
For positive moduli $R$ and a list $B$ of positive moduli,
\[
  \textstyle\prod \fn{radices}(R) = W(R), \qquad
  \fn{newRadix}(B, m) = \dfrac{m}{\gcd(W(B), m)},
\]
and for $m \ge 1$, $\fn{IrreducibleAfter}(B, m) \iff \gcd(W(B), m) = 1$.
\end{theorem}
\leanref{\lean{prod_radices}, \lean{newRadix_eq}, \lean{irreducibleAfter_iff}}

So the product of the radices of positive moduli is their least common multiple, in any order
of the moduli, and a modulus receives its whole value as radix exactly when it is coprime to
the lcm of the moduli before it.

\subsection*{T4 Consecutive moduli}
\begin{theorem}
For $n \ge 2$,
\[
  \fn{newRadix}(\fn{initialModuli}(n), n) =
  \begin{cases}
    p & \text{if } n = p^k,\ p \text{ prime},\ k \ge 1, \\
    1 & \text{otherwise,}
  \end{cases}
\]
\[
  \fn{IrreducibleAfter}(\fn{initialModuli}(n), n) \iff n \text{ is prime},
\]
\[
  \log \fn{newRadix}(\fn{initialModuli}(n), n) = \Lambda(n).
\]
\end{theorem}
\leanref{\lean{newRadix_initialModuli}, \lean{irreducibleAfter_initialModuli_iff},
\lean{log_newRadix_initialModuli}}

\subsection*{T5 Comparison}
\begin{theorem}
For every mixed-radix representation $D$ and modulus $c$,
$\fn{evalMod}(D, c) = \fn{eval}(D) \bmod c$. For $b > 1$ and $\gcd(a, b) = 1$,
$a \cdot \fn{invMod}(a, b) \equiv 1 \pmod b$.
\end{theorem}
\leanref{\lean{eval_append}, \lean{evalMod_eq}, \lean{invMod_spec}}

\begin{theorem}
For positive moduli $R$ and every $n$ there is a mixed-radix representation
$D = ((d_i, r_i))_i$ with
\[
  \fn{toMixedRadix}(R, \st{n}{R}) = D, \qquad (r_i)_i = \fn{radices}(R), \qquad d_i < r_i,
  \qquad \fn{eval}(D) = n \bmod W(R).
\]
For positive moduli $R$ and $n, n' < W(R)$,
$\fn{compareRes}(R, \st{n}{R}, \st{n'}{R}) = \fn{compare}(n, n')$.
\end{theorem}
\leanref{\lean{toMixedRadix_spec}, \lean{compareRes_eq}}

\subsection*{T5c Blocks}
\begin{theorem}
For positive moduli $M$, $\fn{blockModulus}(M, \st{n}{M}) = \big(W(M),\ n \bmod W(M)\big)$;
for blocks $P_1, \dots, P_k$ of positive moduli, $\fn{blockModuli}$ gives the moduli
$W(P_1), \dots, W(P_k)$ with the residues $n \bmod W(P_i)$. For positive moduli $R$, any
moduli $T$ and $n < W(R)$, $\fn{baseExtend}(R, \st{n}{R}, T) = \st{n}{T}$.
\end{theorem}
\leanref{\lean{blockModulus_spec}, \lean{blockModuli_spec}, \lean{baseExtend_spec}}

\begin{theorem}
For blocks $D_1, \dots, D_a$ and quotient blocks $K_1, \dots, K_b$ of positive moduli, with
$Q = \big(W(D_1), \dots, W(D_a), W(K_1), \dots, W(K_b)\big)$,
\[
  \fn{flatten}(\fn{BlockRep.of}(n, D, K)) = (Q, \st{n}{Q}).
\]
If moreover $n$ is below the lcm of all moduli of $D, K$ and $n'$ below that of all moduli of
$D', K'$ (blocks of positive moduli), then
\[
  \fn{compareBlocks}\big(\fn{BlockRep.of}(n, D, K),\ \fn{BlockRep.of}(n', D', K')\big)
  = \fn{compare}(n, n').
\]
\end{theorem}
\leanref{\lean{flatten_spec}, \lean{compareBlocks_eq}}

\subsection*{T6 Primality}
\begin{theorem}
Let $R$ be positive moduli with $n < W(R)$ and $2 < W(R)$. Every answer $b$ of
$\fn{isPrimeRes}(\fn{fuel}, R, \st{n}{R})$ is true exactly when $n$ is prime, and an answer is
given once
\[
  \fn{fuel} > n \quad\text{or}\quad
  \Big(\fn{fuel} > \lfloor\sqrt n\rfloor \ \text{and}\ \big(\lfloor\sqrt n\rfloor + 1\big)^2 < W(R)\Big).
\]
\end{theorem}
\leanref{\lean{isPrimeRes_sound}, \lean{isPrimeRes_eq}}

\subsection*{T7 Interval atoms}
\begin{theorem}
For $k \ge 1$, $\fn{firstMultipleFrom}(lo, k)$ is the least multiple of $k$ that is $\ge lo$,
and it lies in $[lo, lo + k)$. For all $lo$, $hi$, $k$,
\[
  k \in \fn{intervalAtoms}(lo, hi) \iff k \text{ is prime} \ \wedge\
  \exists a:\ lo \le a,\ k \mid a,\ a + k \le hi .
\]
So every atom is prime (\lean{prime_of_mem_intervalAtoms}), and a prime $p$ with
$lo + 2p \le hi + 1$ is an atom: an interval of $L$ consecutive integers yields every prime $p$
with $2p \le L$.
\end{theorem}
\leanref{\lean{firstMultipleFrom_spec}, \lean{mem_intervalAtoms}, \lean{prime_mem_intervalAtoms}}

\begin{theorem}
Every divisor $d \ge 2$ of an element of $\fn{Diff}(lo, hi)$ is in $\fn{Diff}(lo, hi)$, and the
atoms are the minimal elements of $\fn{Diff}(lo, hi)$ under divisibility:
\[
  k \in \fn{intervalAtoms}(lo, hi) \iff k \in \fn{Diff}(lo, hi) \ \wedge\
  \forall d \in \fn{Diff}(lo, hi):\ d \mid k \implies d = k .
\]
\end{theorem}
\leanref{\lean{intervalDifference_of_dvd}, \lean{mem_intervalAtoms_iff_minimal}}

\begin{example}
Checked by the Lean kernel: $\fn{intervalAtoms}(100, 130) = (2, 3, 5, 7, 11, 13, 17)$ and
$\fn{intervalAtoms}(2, 60) = (2, 3, 5, 7, 11, 13, 17, 19, 23, 29)$.
\end{example}

\subsection*{T8 Arithmetic}
\begin{theorem}
For any moduli $R$ and all $x, y, e$:
\[
  \fn{addRes}(R, \st{x}{R}, \st{y}{R}) = \st{x + y}{R}, \quad
  \fn{mulRes}(R, \st{x}{R}, \st{y}{R}) = \st{x y}{R}, \quad
  \fn{powRes}(R, \st{x}{R}, e) = \st{x^e}{R}.
\]
\end{theorem}
\leanref{\lean{addRes_spec}, \lean{mulRes_spec}, \lean{powRes_spec}}

In the rest of this subsection $R$ is a list of positive moduli and $W = W(R)$.

\begin{theorem}
If $y \le x$ then $\fn{subRes}(R, \st{x}{R}, \st{y}{R}) = \st{x - y}{R}$; if $y \le W$ then
$\fn{subRes}(R, \st{x}{R}, \st{y}{R}) = \st{x + (W - y)}{R}$. For $x, y < W$,
\[
  \fn{addOverflow}(R, \st{x}{R}, \st{y}{R}) \iff W \le x + y, \qquad
  \fn{subUnderflow}(R, \st{x}{R}, \st{y}{R}) \iff x < y .
\]
\end{theorem}
\leanref{\lean{subRes_spec_of_le}, \lean{subRes_spec}, \lean{addOverflow_iff},
\lean{subUnderflow_iff}}

\begin{theorem}
For $t \ge 1$ with $t \mid x$, $\fn{exactDivRes}(R, \st{x}{R}, t) = \st{x / t}{R/t}$, where
$R/t = \big(r_i / \gcd(r_i, t)\big)_i$; and if $x < W$ then $x / t < W(R/t)$.
\end{theorem}
\leanref{\lean{exactDivRes_spec}, \lean{div_lt_divModuli}}

\begin{theorem}
For $T \ge 1$ and $x < W$,
$\fn{divModRes}(R, \st{x}{R}, T) = \big(x \bmod T,\ \st{\lfloor x / T \rfloor}{R/T}\big)$. For
$0 < x < W$, with $i = (x - 1) \bmod T + 1$ and $k = \lfloor (x - 1)/T \rfloor + 1$, so that
$x = T (k - 1) + i$ and $1 \le i \le T$,
\[
  \fn{addressRes}(R, \st{x}{R}, T) = \big(i,\ \st{k}{R/T}\big).
\]
\end{theorem}
\leanref{\lean{divModRes_spec}, \lean{addressRes_spec}}

\begin{theorem}
For any moduli $R$, a prime $p$ with $p^v \mid W(R)$ and $p^{v+1} \nmid W(R)$, and every $j$,
\[
  W\big(R, p^{v+1}, p^{v+2}, \dots, p^{v+j}\big) = p^j\, W(R).
\]
\end{theorem}
\leanref{\lean{lcm_append_prime_powers}}

\section{The formalization}\label{sec:formalization}

The formalization is a Lean~4~\cite{lean4} project built against Mathlib~\cite{mathlib}, both
at version \lean{v4.35.0}.
All definitions of Section~\ref{sec:definitions} except the sech comb are in
\lean{lean/Irreducibility/Defs.lean}, which does not import Mathlib; the sech comb needs real
numbers and is defined in \lean{SechComb.lean}. Each group of results has its own file:

\begin{center}
\begin{tabular}{ll@{\qquad}ll}
  T1 Counting & \lean{Residues.lean} & T5 Comparison & \lean{Compare.lean} \\
  T1w Sech comb & \lean{SechComb.lean} & T5c Blocks & \lean{Blocks.lean} \\
  T2 Euclid & \lean{Euclid.lean} & T6 Primality & \lean{Primality.lean} \\
  T3 Mixed-radix system & \lean{Radices.lean} & T7 Interval atoms & \lean{Atoms.lean} \\
  T4 Consecutive moduli & \lean{InitialModuli.lean} & T8 Arithmetic & \lean{Arithmetic.lean}
\end{tabular}
\end{center}

\noindent
\lean{lake build} completes with no \lean{sorry}, \lean{admit} or \lean{native_decide}, and the
theorems depend only on the axioms \lean{propext}, \lean{Quot.sound} and
\lean{Classical.choice}. The examples quoted above are checked by the kernel with
\lean{decide}. To build:
\begin{verbatim}
cd lean
lake exe cache get
lake build
\end{verbatim}

\section{Use of AI}

The construction (residues computed by counting, the mixed-radix system of arbitrary moduli
built modulus by modulus, block representations, the residue read from a sech comb, the
irreducibility of a modulus after other moduli, the atoms of an interval, arithmetic on
residues) and the choice of the statements to be proved are the author's. Generative AI
(Claude Opus 5.5, Anthropic, used through Claude Code) contributed in the following ways:
\begin{itemize}
  \item It wrote the Lean definitions from the author's description of the construction, and
    the Lean proofs. The author checked the definitions against the construction; the
    correctness of the proofs rests on the Lean kernel.
  \item It wrote the statements in standard notation and drafted this text. The author set the
    structure of the paper and the stated contributions.
  \item It searched the literature and checked the identifiers of the references. The author
    decided how the work is placed relative to that literature.
\end{itemize}

\begin{thebibliography}{9}

\bibitem{didier}
L.-S.~Didier, L.~Glandus, N.~El~Mrabet and J.-M.~Robert,
\emph{Residue number system comparison revisited, a software perspective},
2025 5th International Conference on Electrical, Computer and Energy Technologies (ICECET),
2025; \href{https://arxiv.org/abs/2605.18415}{arXiv:2605.18415}.

\bibitem{eberl-afp}
M.~Eberl,
\emph{Elementary facts about the distribution of primes},
Archive of Formal Proofs, February 2019, ISSN 2150-914x,
\url{https://isa-afp.org/entries/Prime_Distribution_Elementary.html}.

\bibitem{garner}
H.~L.~Garner,
\emph{The residue number system},
IRE Transactions on Electronic Computers EC-8 (1959), no.~2, 140--147,
\href{https://doi.org/10.1109/TEC.1959.5219515}{doi:10.1109/TEC.1959.5219515}.

\bibitem{mathlib}
The mathlib Community,
\emph{The Lean mathematical library},
Proceedings of the 9th ACM SIGPLAN International Conference on Certified Programs and Proofs
(CPP 2020), 367--381,
\href{https://doi.org/10.1145/3372885.3373824}{doi:10.1145/3372885.3373824}.

\bibitem{lean4}
L.~de~Moura and S.~Ullrich,
\emph{The Lean 4 theorem prover and programming language},
Automated Deduction -- CADE 28, Lecture Notes in Computer Science, 2021, 625--635,
\href{https://doi.org/10.1007/978-3-030-79876-5_37}{doi:10.1007/978-3-030-79876-5\_37}.

\bibitem{sundaram}
S.~Sundaram and C.~N.~Hadjicostis,
\emph{Fault-tolerant convolution via Chinese remainder codes constructed from non-coprime
moduli},
IEEE Transactions on Signal Processing 56 (2008), no.~9, 4244--4254,
\href{https://doi.org/10.1109/TSP.2008.924636}{doi:10.1109/TSP.2008.924636}.

\end{thebibliography}

\end{document}
